\chapter{Conclusioni e sviluppi futuri}
\label{chap:conclusioni}

\section{Bilancio del lavoro}

Il lavoro ha esaminato come quattro approcci al rilevamento delle anomalie visive industriali descrivano la normalità e riconoscano i difetti. Il confronto tra CCW-SSIM, PaDiM, RD4AD e FastFlow ha collegato i principi dello stato dell'arte alle configurazioni adottate, utilizzando le immagini della categoria \texttt{leather} di MVTec AD e un protocollo comune di valutazione. Le metriche, le mappe di anomalia e i casi di errore hanno permesso di considerare insieme il riconoscimento delle immagini difettose e la localizzazione delle regioni anomale.

PaDiM, RD4AD e FastFlow hanno separato completamente le immagini normali e difettose del test esaminato. La localizzazione ha invece mostrato differenze tra i tre modelli: FastFlow ha ottenuto i valori più elevati delle metriche considerate, seguito da RD4AD e PaDiM. Le mappe hanno reso visibile che una risposta nella posizione corretta può estendersi oltre il contorno annotato; precisione e richiamo hanno descritto il compromesso tra copertura del difetto e segnalazione di pixel normali. Il risultato a livello di immagine va quindi letto insieme alla qualità della localizzazione.

CCW-SSIM ha presentato una separazione delle immagini meno efficace e una copertura regionale inferiore. Negli esempi analizzati, le risposte alla texture normale rendono meno netto il contrasto del difetto e possono contribuire a una segnalazione anche in assenza di anomalie. La lettura dei casi visivi ha così aiutato a comprendere aspetti del comportamento della pipeline che la sola AUROC non descrive pienamente.

Anche i tempi di apprendimento riflettono le differenze tra gli approcci. PaDiM ha richiesto il minor tempo di adattamento locale, poiché stima le statistiche delle feature senza aggiornare l'estrattore. CCW-SSIM ha richiesto il tempo maggiore, mentre RD4AD e FastFlow si sono collocati in posizione intermedia. Le misure descrivono la preparazione delle pipeline nelle condizioni adottate; il costo del pre-addestramento esterno e le risorse necessarie durante l'inferenza richiedono valutazioni distinte.

Il confronto con le pubblicazioni ha mostrato AUROC di PaDiM, RD4AD e FastFlow prossime ai valori riportati per \texttt{leather}. Le differenze di configurazione e valutazione aiutano però a definire la portata di questo riscontro. Per RD4AD cambia il protocollo di calcolo del PRO; per CCW-SSIM differiscono la metrica della fonte e la formulazione della componente cromatica. Le verifiche diagnostiche hanno approfondito alcuni aspetti della variante locale, senza isolare il beneficio del colore rispetto a CW-SSIM. La valutazione di tale contributo rimane quindi una domanda aperta.

Il contributo dell'elaborato consiste in un confronto documentato che mette in relazione prestazioni, costi di adattamento e comportamento delle mappe. La registrazione di dati, configurazioni e risultati permette di ricondurre le osservazioni alle esecuzioni analizzate. Nel caso di studio, questa impostazione ha consentito di riconoscere sia le capacità delle diverse pipeline sia i punti che richiedono un approfondimento sperimentale.

\section{Portata delle conclusioni}

I risultati riguardano un unico materiale e le configurazioni utilizzate nel progetto. Le cinque tipologie di difetto permettono di esaminare manifestazioni diverse delle anomalie all'interno di \texttt{leather}, mentre la generalizzazione ad altri materiali richiede ulteriori esperimenti. Le differenze di risoluzione, pre-addestramento e trattamento delle mappe sono parte del confronto e non sono state isolate mediante prove controllate.

Una sola esecuzione per modello non permette inoltre di quantificare la variabilità dei risultati. I piccoli scarti tra pipeline descrivono le esecuzioni osservate, ma non dimostrano una superiorità statisticamente significativa. Anche la scelta delle soglie delimita il significato delle prestazioni: il massimo F1 è calcolato sul test stesso e la proporzione di immagini difettose può differire da quella di una linea produttiva. Il passaggio all'impiego industriale richiederebbe soglie selezionate su dati separati e una valutazione del compromesso tra omissioni e falsi allarmi in funzione dell'applicazione.

Queste condizioni definiscono anche le priorità per lo sviluppo del lavoro: approfondire i componenti delle pipeline, verificare la stabilità delle prestazioni e valutare il comportamento in condizioni operative differenti.

\section{Sviluppi futuri}

Un primo approfondimento riguarda il contributo dell'informazione cromatica in CCW-SSIM. La diversa pesatura tra addestramento e inferenza motiva una revisione controllata della componente cromatica, con parametri coerenti con la normalizzazione. Il confronto con CW-SSIM, mantenendo invariati architettura, dati, addestramento e valutatore, permetterebbe poi di misurare il beneficio del colore. L'analisi della tipologia \texttt{color} andrebbe affiancata a quella degli altri difetti e delle risposte sulla texture normale. La modifica della perdita richiederebbe un nuovo addestramento e risultati identificati come un esperimento successivo.

La stabilità del confronto potrebbe essere esaminata attraverso repliche con diversi semi pseudo-casuali. All'interno di \texttt{leather}, il calcolo delle metriche di localizzazione per tipologia di difetto completerebbe inoltre i conteggi di classificazione già disponibili. Il confronto con i valutatori originali, soprattutto per il PRO di RD4AD, aiuterebbe a distinguere le differenze di calcolo da quelle delle pipeline.

Un altro sviluppo riguarda la scelta del punto operativo. Un insieme di validazione separato consentirebbe di definire le soglie prima della valutazione sul test. Per CCW-SSIM, il confronto tra il massimo della mappa e aggregazioni alternative permetterebbe di verificare se una minore sensibilità ai picchi isolati riduca i falsi positivi mantenendo il rilevamento dei piccoli difetti.

La valutazione applicativa richiederebbe infine misure di latenza, numero di immagini elaborate per unità di tempo e memoria di inferenza sullo stesso hardware, distinguendo preprocessing, esecuzione del modello e post-elaborazione. Prove controllate con variazioni di illuminazione e colore permetterebbero di esaminare la robustezza alle condizioni di acquisizione; l'estensione ad altri materiali consentirebbe di studiare la generalizzazione.

Il caso di studio mostra quindi l'utilità di affiancare alla capacità di riconoscere un'immagine difettosa l'analisi della localizzazione, delle soglie e dei costi osservati. I metodi basati su feature hanno ottenuto i risultati più favorevoli nelle configurazioni adottate. Le domande emerse dal confronto offrono una base per verificarne la stabilità e approfondire il comportamento dei modelli in nuovi esperimenti.
